%%%%%%%%%%%%%%%%%%%%%%%%%%%%%%%%%%%%%%%%%%%%%%%%%%%%%%%%%%%%%%%%%%%%%%
% Lớp 11 | Đường thẳng, mặt phẳng. Quan hệ song song trong | Hai mặt phẳng song song
% Dạng 1H4?4-5: Xác định mặt phẳng chứa đường thẳng (hoặc đi qua hai điểm) và song
% Số câu: 21  (sinh tự động bởi tach_cau.py; từ nay thầy sửa trực tiếp file này)
%%%%%%%%%%%%%%%%%%%%%%%%%%%%%%%%%%%%%%%%%%%%%%%%%%%%%%%%%%%%%%%%%%%%%%

%%%%%%%%%%%%%%%% TRẮC NGHIỆM 4 PHƯƠNG ÁN %%%%%%%%%%%%%%%%

% nguồn: Ok_Toan11_TracNghiem_Lan2
\begin{ex}%[1H4H4-5]%[ma:637f1f3d]
	Cho  hình hộp $ABCD.A'B'C'D'$. Mặt phẳng $\left(AB'D'\right)$ song song với mặt phẳng nào sau đây?
	 \choice
		 {\True $\left(BDC'\right)$}
		{$\left(BCA'\right)$} 
		{$\left(BDA'\right)$} 
		 {$\left(A'C'C\right)$}
	\loigiai{ 
		\begin{center}
			\begin{tikzpicture}[scale=0.8,font=\footnotesize,line join=round,line cap=round,>=stealth]
					\coordinate (B) at  (0,0);
					 \coordinate (A) at (-2,-2);
					\coordinate (C) at (6,0); 
					 \coordinate (D) at ($(A)+(C)-(B)$);
					\foreach \i in {A,B,C,D}{\coordinate  (\i') at ($(\i)-(0,4)$);}
					\draw  (A)--(B)--(C)--(D)--cycle;
					\draw  (A)--(A')
						(C)--(C') 
						 (D)--(D')
						 (A')--(D')--(C')
						(B)--(D)--(C') 
						(A)--(D'); 
					 \draw[dashed,thin] (B)--(B')--(A')
						(B')--(C')--(B) 
						 (A)--(B')
						 (B')--(D');
					 \foreach \i/\g in {A'/170,B'/170,C'/-90,D'/-90,A/90,B/90,C/-20,D/10}{\draw[fill=blue](\i) circle (1pt)
						 ($(\i)+(\g:5mm)$) node[scale=1]{$\i$};}
				\end{tikzpicture}
		\end{center}
		Mặt phẳng $\left(AB'D'\right)$  song song với mặt phẳng $\left(BDC'\right)$.\\
		Thật vậy, ta có $AB' \parallel DC'$  và $AD' \parallel BC'$ có điều cần chứng minh.
	} 
\end{ex}

% nguồn: Ok_Toan11_TracNghiem_Lan2
\begin{ex}%[1H4H4-5]%[ma:0f0b4906]
	Cho hình hộp $ABCD. A' B'  C' D'$. Gọi $M$ là trung điểm của $AB$. Mặt phẳng $\left(MA' C'\right)$ cắt hình hộp $ABCD. A' B' C' D'$ theo thiết diện là hình gì?
	 \choice
		{\True Hình  thang}
		{Hình ngũ giác} 
		{Hình  lục giác}
		{Hình tam  giác}
	 \loigiai{
		\begin{center}
			\begin{tikzpicture}[scale=.8,font=\footnotesize,line join=round,line cap=round,>=stealth]
					\coordinate[label=above left:$A$]  (A) at (0,0);
					\coordinate[label=below:$C$] (C) at  (2,-1.5);
					 \coordinate[label=right:$D$] (D) at (4,0);
					\coordinate[label=below left:$B$]  (B) at (-2,-1.5);
					 \coordinate[label=above:$A'$] (A') at (1,4);
					\coordinate[label=below right:$C'$] (C') at (3,2.5); 
					\coordinate[label=above:$D'$]  (D') at (5,4);
					\coordinate[label=left:$B'$]  (B') at (-1,2.5);
					\coordinate[label=above:$O$] (O) at ($(A')!0.5!(C')$); 
					 \coordinate[label=above:$N$] (N) at ($(B)!0.5!(C)$);
					\coordinate[label=left:$M$] (M)  at ($(A)!0.5!(B)$);
					\path (intersection of C'--N and A'--M) coordinate (I) node[below]{$I$}; 
					 \draw (D')--(A')--(B')--(C')--(D')--(D)--(C)--(C')
						(C)--(B)--(B') 
						(A')--(C') --(I)--(B) 
						 (B')--(D');
					 \draw[dashed] (B)--(A)--(D)
						(A)--(A')--(I); 
					\foreach  \diem in {A,B,C,D,A',B',C',D',O,M,N,I}\fill (\diem) circle (1pt);
				\end{tikzpicture}
		\end{center}
		Trong mặt phẳng $\left(ABB' A'\right), AM$  cắt $BB'$ tại $I$.\\
		Do $MB\parallel A' B'; MB=\dfrac{1}{2}  A' B'$ nên $B$ là trung điểm $B' I$ và $M$ là trung điểm của $IA'$.\\
		Gọi $N$ là giao  điểm của $BC$ và $C' I$.\\
		Do $BN\parallel B' C$ và $B$ là trung điểm $B' I$  nên $N$ là trung điểm của $C' I$.\\
		Suy ra tam giác $IA' C'$ có $MN$  là đường trung bình.\\
		Ta có mặt phẳng $\left(MA' C'\right)$ cắt hình hộp  $ABCD. A' B' C' D'$ theo thiết diện là tứ giác $A' MNC'$ có $MN\parallel A' C'$.\\
		Vậy thiết  diện là hình thang $A' MNC'$.
	} 
\end{ex}

% nguồn: Ok_Toan11_TracNghiem_Lan2
\begin{ex}%[1H4H4-5]%[ma:174dbe8e]
	Cho tứ diện đều $SABC$. Gọi $I$ là trung điểm của đoạn $AB$, $M$ là điểm di động trên đoạn $AI$. Qua $M$ vẽ mặt phẳng  $(\alpha)$ song song với $(SIC)$. Thiết diện tạo bởi $(\alpha)$ với tứ diện $SABC$ là
	 \choice
		{hình  bình hành}
		{\True tam  giác cân tại $M$}
		{tam giác  đều}
		{hình  thoi}
	 \loigiai{
		\begin{center}
			\begin{tikzpicture}[scale=1,font=\footnotesize,line join=round,line cap=round,>=stealth]
					 \coordinate[label=above:$S$] (S) at (2.5,3.5);
					\coordinate[label=below:$B$] (B)  at (3,-1.5);
					\coordinate[label=right:$C$] (C) at  (5,0);
					\coordinate[label=left:$A$] (A) at  (0,0);
					 \coordinate[label=below:$I$] (I) at ($(A)!0.5!(B)$);
					\coordinate[label=below:$M$] (M) at ($(A)!0.5!(I)$); 
					\coordinate[label=below:$P$] (P) at  ($(A)!0.5!(C)$);
					 \coordinate[label=left:$N$] (N) at ($(A)!0.5!(S)$);
					\draw (S)--(A)--(B)--(C)--(S)--(B) 
						(S)--(I) 
						(M)--(N); 
					 \draw[dashed] (A)--(C)
						(N)--(P)--(M) 
						 (C)--(I);
					\foreach  \diem in {A,B,C,S,I,M,P,N}\fill (\diem) circle (1pt);
				\end{tikzpicture}
		\end{center}
		Qua $M$  vẽ $MP\parallel IC, P\in AC, MN\parallel SI, N\in SA$.\\
		Do đó thiết diện tạo  bởi $(\alpha)$ với tứ diện $SABC$ là tam giác $MNP$.\\
		Ta có  $\dfrac{MN}{SI}=\dfrac{MP}{IC}$ và $SI=IC$, suy ra $MN=MP$.\\
		Vậy thiết diện là tam giác  $MNP$ cân tại $M$.
	} 
\end{ex}

% nguồn: Ok_Toan11_TracNghiem_Lan2
\begin{ex}%[1H4H4-5]%[ma:65987b7a]
	Cho hình chóp $S. ABCD$ có đáy $ABCD$ là hình bình hành tâm $O$ có $AC=a, BD=b$. Tam giác $SBD$ là tam giác đều. Một mặt phẳng $(\alpha)$ di động song song với mặt phẳng $(SBD)$ và đi qua điểm $I$ trên đoạn $AC$ và  $AI=x \,(0< x < a)$. Thiết diện của hình chóp cắt bởi $(\alpha)$ là hình gì?
	\choice 
		 {Hình bình hành}
		 {\True Tam giác}
		 {Tứ giác}
		 {Hình thang}
	\loigiai{ 
		\begin{center}
			\begin{tikzpicture}[scale=1,font=\footnotesize,line join=round,line cap=round,>=stealth]
					 \coordinate[label=above left:$A$] (A) at (0,0);
					\coordinate[label=below:$C$] (C)  at (2,-1.25);
					\coordinate[label=right:$D$]  (D) at (4,0);
					\coordinate[label=below:$B$] (B)  at (-2,-1.25);
					\coordinate[label=above:$S$] (S) at  (1,4);
					\coordinate[label=below:$O$] (O) at  ($(A)!0.5!(C)$);
					\coordinate[label=below:$I$]  (I) at ($(A)!0.55!(O)$);
					\coordinate[label=above  right:$N$] (N) at ($(A)!0.55!(D)$);
					 \coordinate[label=left:$M$] (M) at ($(A)!0.55!(B)$);
					\coordinate[label=left:$P$] (P) at ($(A)!0.55!(S)$); 
					\draw (S)--(B)--(C)--(S)--(D)--(C); 
					\draw[dashed]  (B)--(A)--(D)--(B)
						 (S)--(A)--(C)
						(M)--(N)--(P)--(M); 
					\foreach \diem in {A,B,C,D,S,O,I,M,N,P}\fill  (\diem) circle (1pt);
				\end{tikzpicture}
		\end{center}
		\textbf{Trường hợp 1.} Xét $I$  thuộc đoạn $OA$.\\
		Ta có $\heva{&I\in(\alpha) \cap(ABD)\\&(\alpha) \parallel (SBD) \\&(ABD) \cap(SBD)=BD}  \Rightarrow(\alpha) \cap(ABD)=MN\parallel BD, I\in MN$.\\
		Tương tự  $\heva{&N\in(\alpha) \cap(SAD)\\
		 &(\alpha) \parallel (SBD)\\
		 &(SAD)\cap(SBD)=SD
		} \Rightarrow(SAD)\cap(\alpha)=NP\parallel  SD, P\in SN$.\\
		Thiết diện là tam giác  $MNP$.\\
		Do $ \heva{&(\alpha) \parallel  (SBD)\\
		 &(SAB) \cap(SBD)=SB \\
		 &(SAB) \cap(\alpha)=MP
		}\Rightarrow  MP\parallel SB$.\\
		Ta  có $\triangle MNP\sim \triangle BDS$, mà tam giác $BDS$ đều nên tam giác $MNP$ đều.\\
		\textbf{Trường hợp 2.} Điểm $I$ thuộc đoạn $OC$, tương tự trường hợp 1 ta  được thiết diện là tam giác đều $HKL$.
	} 
\end{ex}

% nguồn: Ok_Toan11_TracNghiem_Lan2
\begin{ex}%[1H4H4-5]%[ma:a510bbae]
	Cho hình chóp $S. ABCD$ có đáy $ABCD$ là hình bình hành và $M, N$ lần lượt là trung điểm của $AB, CD$. Xác định thiết diện của hình chóp cắt bởi $(\alpha)$ đi qua $MN$ và song song với mặt phẳng  $(SAD)$.Thiết diện là hình gì?
	 \choice
		{\True Hình thang} 
		{Hình  bình hành}
		 {Tứ giác}
		{Tam giác} 
	 \loigiai{
		\begin{center}
			\begin{tikzpicture}[scale=1,font=\footnotesize,line join=round,line cap=round,>=stealth]
					\coordinate[label=above left:$A$]  (A) at (0,0);
					 \coordinate[label=below:$C$] (C) at (2,-1.25);
					\coordinate[label=right:$D$] (D) at  (4,0);
					\coordinate[label=below:$B$] (B) at (-2,-1.25); 
					\coordinate[label=above:$S$]  (S) at (1,4);
					 \coordinate[label=left:$M$] (M) at ($(A)!0.5!(B)$);
					\coordinate[label=left:$K$] (Q) at  ($(B)!0.5!(S)$);
					 \coordinate[label=right:$N$] (N) at ($(D)!0.5!(C)$);
					\coordinate[label=right:$H$] (P) at ($(C)!0.5!(S)$); 
					\draw  (S)--(B)--(C)--(S)--(D)--(C)
						(N)--(P)--(Q); 
					 \draw[dashed] (B)--(A)--(D)
						 (S)--(A)
						(M)--(N) 
						 (M)--(Q);
					\foreach \diem in {A,B,C,D,S,N,M,Q,P}\fill (\diem) circle (1pt); 
				\end{tikzpicture}
		\end{center}
		Ta có $\heva{&M\in(SAB) \cap(\alpha) \\& (SAB)  \cap(SAD)=SA} \Rightarrow(SAB) \cap(\alpha)=MK\parallel SA, K\in SB$.\\
		Tương tự $\heva{&N\in(SCD) \cap(\alpha)\\&(\alpha) \parallel (SAD)\\&(SCD) \cap(SAD)=SD} \Rightarrow(SCD) \cap(\alpha)=NH\parallel SD, H\in  SC$.\\
		Dễ thấy $HK=(\alpha) \cap(SBC)$. Thiết diện  là tứ giác $MNHK$.\\
		Ba mặt phẳng $(ABCD),(SBC)$ và $(\alpha)$ đôi một cắt nhau theo các giao  tuyến là $MN, HK, BC$, mà $MN\parallel BC$, suy ra $ MN\parallel HK$. Vậy thiết diện là một hình thang.
	 }
\end{ex}

% nguồn: Ok_Toan11_TracNghiem_Lan2
\begin{ex}%[1H4H4-5]%[ma:38dc7a4c]
	Cho hình vuông $ABCD$ và tam giác đều $SAB$ nằm trong  hai mặt phẳng khác nhau. Gọi $M$ là điểm di động trên đoạn $AB$. Qua $M$ vẽ mặt phẳng $(\alpha)$ song song với $(SBC)$. Thiết diện tạo bởi $(\alpha)$ và hình chóp $S. ABCD$ là hình gì?
	 \choice
		{Hình  tam giác}
		{Hình bình  hành}
		{Hình thang} 
		 {\True Hình vuông}
	 \loigiai{
		\begin{center}
			\begin{tikzpicture}[scale=1,font=\footnotesize,line join=round,line cap=round,>=stealth]
					\coordinate[label=above left:$A$] (A) at (0,0); 
					\coordinate[label=below:$C$] (C) at  (2,-1.25);
					\coordinate[label=right:$D$] (D) at  (4,0);
					\coordinate[label=below:$B$] (B)  at (-2,-1.25);
					 \coordinate[label=above:$S$] (S) at (1,4);
					 \coordinate[label=above:$O$] (O) at ($(A)!0.5!(C)$);
					\coordinate[label=left:$M$] (M)  at ($(A)!0.6!(B)$);
					\coordinate[label=left:$Q$]  (Q) at ($(A)!0.6!(S)$);
					\coordinate[label=right:$N$] (N) at ($(D)!0.6!(C)$); 
					 \coordinate[label=right:$P$] (P) at ($(D)!0.6!(S)$);
					\draw  (S)--(B)--(C)--(S)--(D)--(C)
						(N)--(P); 
					\draw[dashed] (B)--(A)--(D)--(B) 
						(S)--(A)--(C) 
						(M)--(N) 
						 (M)--(Q)--(P);
					\foreach \diem in {A,B,C,D,S,O,N,M,Q,P}\fill (\diem) circle (1pt); 
				\end{tikzpicture}
		\end{center}
		Lần lượt lấy các điểm $N, P, Q$ thuộc các cạnh $CD, SD, SA$ thỏa $MN\parallel BC, NP\parallel SC, PQ\parallel AD$.  \\
		Suy ra $(\alpha) \equiv(MNPQ)$ và $(\alpha)  \parallel(SBC)$.\\
		Theo cách dựng trên thì thiết diện là hình thang. 
	} 
\end{ex}

% nguồn: Ok_Toan11_TracNghiem_Lan2
\begin{ex}%[1H4V4-5]%[ma:445c3542]
	Cho hình chóp $S.ABCD$ có đáy $ABCD$ là hình thang cân với cạnh bên $BC=3$. Hai  đáy $AB=8$ và $CD=4$. Mặt phẳng $(P)$ song song với mặt phẳng $(ABCD)$ và cắt cạnh $SA$ tại $M$ sao cho $SA=3SM$. Diện tích thiết diện của $(P)$ và hình chóp $S.ABCD$ bằng bao nhiêu?
	 \choice
		{\True  $\dfrac{2 \sqrt{5}}{9}$}
		{$\dfrac{2 \sqrt{5}}{3}$} 
		{$\dfrac{7 \sqrt{3}}{3}$} 
		 {$\dfrac{7 \sqrt{3}}{9}$}
	\loigiai{ 
		\begin{center}
			\begin{tikzpicture}[scale=1,font=\footnotesize,line join=round,line cap=round,>=stealth]
					\coordinate (S)  at (2,4);
					\coordinate (A) at  (0,0);
					 \coordinate (D) at (1,-2);
					\coordinate (C) at (3,-2); 
					\coordinate (B)  at (5,0);
					\coordinate (M) at  ($(S)!0.333!(A)$);
					\coordinate  (N) at ($(S)!0.333!(B)$);
					\coordinate (P) at ($(S)!0.333!(C)$); 
					 \coordinate (Q) at ($(S)!0.333!(D)$);
					\coordinate (H) at ($(M)!0.5!(N)$); 
					\fill[yellow](M)--(N)--(P)--(Q)--cycle; 
					 \draw (S)--(A)
						(S)--(B) 
						(S)--(C) 
						 (S)--(D)
						(B)--(C) 
						(A)--(D)--(C)--(B) 
						 (M)--(N)--(P)--(Q)
						 (M)--(Q);
					 \draw[dashed,thin](A)--(B)
						(P)--(H); 
					\foreach \i/\g  in {S/90,A/180,B/-90,C/-90,D/-90,M/180,Q/-45,P/-45,N/45,H/135}{\draw[fill=black](\i) circle (1pt)
						 ($(\i)+(\g:3mm)$) node[scale=1]{$\i$};}
				\end{tikzpicture}
		\end{center}
		Gọi $M$, $N$, $P$ lần lượt là giao điểm của $(P)$ với các cạnh $SB$, $SC$, $SD$. Suy ra thiết diện của $(P)$  với hình chóp là tứ giác $MNPQ$.\\
		Do $(P)\parallel(ABCD)$ nên  $\dfrac{MN}{AB}=\dfrac{NP}{BC}=\dfrac{PQ}{CD}=\dfrac{SM}{SA}=\dfrac{1}{3}$\\
		$\Rightarrow MN=\dfrac{8}{3}$, $NP=1$, $P Q=\dfrac{4}{3}$,  $QM=1$.\\
		Gọi $H$ là hình chiếu của  $P$ trên cạnh $M N$.\\
		Do $MNPQ$ là hình thang  cân nên \[HN=\dfrac{MN-PQ}{2}=\dfrac{2}{3} \Rightarrow PH=\sqrt{PN^2-HN^2}=\dfrac{\sqrt{5}}{3}.\]
		Vậy  $S_{MNPQ}=\dfrac{1}{2} \cdot P H \cdot(M N+PQ)=\dfrac{2 \sqrt{5}}{3}$.
	} 
\end{ex}

% nguồn: Ok_TracNghiemLop11_Lan1
\begin{ex}%[1H4V4-5]%[ma:a834951c]
	Cho hình chóp $S.ABCD$ có đáy $ABCD$ là hình bình hành, mặt bên $SBC$ là tam giác đều. Gọi $M$ là điểm di động trên đoạn thẳng $AB$, $M\ne A;M\ne  B$. Qua $M$ dựng mặt phẳng $(\alpha)$ song song với mặt phẳng $(SBC)$. Thiết diện tạo với mặt phẳng $(\alpha)$ và hình chóp $S.ABCD$ là hình gì?
	\choice 
		 {\True Hình thang cân}
		{Hình  thang vuông}
		{Hình  tam giác}
		{Hình  bình hành}
	\loigiai{ 
	\immini{Ta có $(ABCD)\cap (SBC)=BC$, $M\in (\alpha)\cap (ABCD)$ và $(\alpha)\parallel (SBC)$ nên giao tuyến của $(\alpha)$ và $(ABCD)$  là đường thẳng đi qua $M$ và song song với $BC$, cắt $DC$ tại $N$.\\
		Tương  tự,
		\begin{itemize}
			\item giao tuyến của $(\alpha)$ và $(SAB)$ là đường thẳng đi qua $M$ và  song song với $SB$, cắt $SA$ tại $Q$;
			\item giao tuyến của $(\alpha)$ và $(SCD)$ là đường chứa $N$ và song song với $SC$, cắt $SD$  tại $P$.
		\end{itemize}
		}{\begin{tikzpicture}[scale=1,font=\footnotesize,line join=round,line cap=round,>=stealth]
		\def\k{0.6}; 
		\path (0,0) coordinate (D)++(3,0)  coordinate (C)++(1.1,1) coordinate (B)
			 ($(B)+(D)-(C)$) coordinate (A)
			(0.7,2.5) coordinate (h) 
			 ($(A)+(h)$) coordinate (S)
			($(A)!{\k}!(B)$) coordinate  (M)
			($(D)!{\k}!(C)$)  coordinate (N)
			 ($(A)!{\k}!(S)$) coordinate (Q)
			 ($(D)!{\k}!(S)$) coordinate (P);
		\draw (S)--(D)--(C)--(B)--cycle  (S)--(C)
			(P)--(N); 
		 \draw[dashed] (S)--(A)--(D)
			 (A)--(B)
			 (P)--(Q)--(M)--(N);
		\foreach \p/\g in {A/150,D/240,C/-60,B/0,S/90,M/45,N/-90,P/120,Q/0}  \fill[black] (\p) circle (1pt)+(\g:0.3) node{$\p$};
	\end{tikzpicture}
		}\noindent\\ 
		Khi đó, thiết diện là hình thang $MNPQ$.\\ 
		Ta có $(\alpha)\cap (SAD)=PQ$, $(\alpha)\cap (ABCD)=MN$, $(SAD)\cap (ABCD)=AD$ và $MN\parallel  BC\parallel AD$ nên $PQ\parallel MN$.\\
		Suy  ra $MNPQ$ là hình thang.\\
		Dễ thấy được  $\dfrac{PN}{SC}=\dfrac{DN}{DC}=\dfrac{AM}{AB}=\dfrac{QM}{SB}$.\\
		Hơn nữa, tam giác $SBC$ đều nên $SB=SC$. Do đó $MQ=NP$.\\ 
		Vậy $MNPQ$ là hình thang  cân.
	} 
\end{ex}


%%%%%%%%%%%%%%%% ĐÚNG/SAI %%%%%%%%%%%%%%%%

% nguồn: Ok_dungsai_lop11_Lan2
\begin{ex}%[1H4H4-5]%[ma:e3d5d5c6]
	Cho hình hộp $ABCD.A'B'C'D'$.  Gọi $M$, $N$, $P$, $Q$ lần lượt là trung điểm của các cạnh $BC$, $AA'$, $C'D'$, $AD'$.
	\choiceTF 
		{\True $NQ \parallel A'D'$ và $NQ=\dfrac{1}{2} A'D'$} 
		{\True Tứ giác $MNQC$ là hình  bình hành}
		{\True $MN \parallel  \left(ACD'\right)$}
		{\True $\left(MNP\right)  \parallel \left(ACD'\right)$}
	 \loigiai{
		\begin{center}
			\begin{tikzpicture}[scale=0.8,font=\footnotesize,line join=round,line cap=round,>=stealth]
					\coordinate (A)  at (0,0);
					\coordinate  (B) at (-2,-2);
					 \coordinate (D) at (6,0);
					\coordinate  (C) at ($(B)+(D)-(A)$);
					\foreach \i in {A,B,C,D}{\coordinate (\i') at  ($(\i)+(0,6)$);}
					\draw  (B)--(C)--(D);
					\draw  (B)--(B')
						(C)--(C') 
						(D)--(D') 
						 (A')--(D')--(C')--(B')--(A');
					\draw[dashed,thin] (A')--(A)--(B) 
						(A)--(D); 
					\coordinate (M)  at ($(B)!0.5!(C)$);
					 \coordinate (N) at ($(A)!0.5!(A')$);
					\coordinate (P) at ($(C')!0.5!(D')$); 
					\coordinate (Q)  at ($(A)!0.5!(D')$);
					\coordinate (O)  at ($(A)!0.5!(C)$);
					\draw  (C)--(D');
					\draw[dashed,thin] (M)--(N)--(Q)--(C); 
					\draw[dashed,thin] (N)--(P)--(M); 
					\draw[dashed,thin]  (A)--(D')--(O)--(M)
						 (A)--(C);
					\foreach \i/\g in {A'/170,B'/170,C'/-10,D'/0,A/-170,B/140,C/-20,D/10,O/-90,M/130,N/180,P/10,Q/90}{\draw[fill=blue](\i) circle  (1pt)
						($(\i)+(\g:5mm)$) node[scale=1]{$\i$};} 
				\end{tikzpicture}
		\end{center}
		\begin{itemchoice}
			\itemch Ta có $NQ$ là đường trung bình  của $\Delta AA'D'$ nên $NQ \parallel A'D'$ và $NQ=\dfrac{1}{2} A'D'$.
			\itemch Ta có $M$ là trung điểm $BC,BC$  lại song song và bằng $A'D'$ nên $MC$ song song và độ dài bằng nửa $A'D'$. Kết hợp với két quả câu (1) suy ra $MC$ song song và bằng $NQ$.\\
			Do đó $MNQC$ là hình  bình hành.
			\itemch Do $MNQC$ là hình bình hành nên $MN$ song song với $QC$  mà $QC$ nằm trong $\left(ACD'\right)$ nên $MN$ song song với $\left(ACD'\right)$.
			\itemch Trước tiên ta sẽ chứng minh $MOD'P$ là hình bình hành. Thật vậy gọi $O$ là trung điểm $AC$ thì suy ra $MO$ là đường trung bình của $\triangle  ABC$.\\
			Do  đó $OM$ song song và có độ dài bằng nửa đoạn $AB$.\\
			Lại có $P$ là trung điểm $C'D'$ mà $C'D'$ song song và bằng $AB$  nên $PD'$ song song và độ dài bằng nửa đoạn $AB$.\\
			Suy ra $OM$ song song và bằng $PD'$ do đó $MOD'P$  là hình bình hành.\\
			Từ đó ta có $MP$ song song với $D'O$  mà $D'O\subset \left(ACD'\right)$ nên $MP$ song song với $\left(ACD'\right)$.\\
			Ta đã chứng minh được $MN$ song song $\left(ACD'\right)$.\\ 
			Mặt khác $MP$, $MN$ nằm trong $\left(MNP\right)$  và $MP$, $MN$ cắt nhau.\\
			Vậy $\left(MNP\right)$ song song với  $\left(ACD'\right)$.
		\end{itemchoice}
	} 
\end{ex}


%%%%%%%%%%%%%%%% TRẢ LỜI NGẮN %%%%%%%%%%%%%%%%

% nguồn: Ok_TraLoiNgan_Lop11_Lan1
\begin{ex}%[1H4H4-5]%[ma:e4cb930b]
	Cho hình chóp $S.ABCD$ có đáy $ABCD$ là hình thang với $AD$ là đáy lớn và $AD=2BC$. Gọi  $M$, $N$ lần lượt là trung điểm của $SA$ và $AD$. Biết mặt phẳng $(\alpha)$ đi qua $MN$ và song song với $(SCD)$ cắt cạnh $BC$ tại điểm $E$. Tính tỉ số $\dfrac{CE}{CB}$.
	\shortans{1} 
	 \loigiai{
	\immini{Ta  có $(\alpha)\cap (ABCD)=NE$, $(SCD)\cap (ABCD)=CD$ và $(\alpha)\parallel (SCD)$ nên $NE\parallel CD$.\\
		 Theo đề bài thì $ND\parallel BC$.\\
		Do đó  $NECD$ là hình bình hành. Suy ra $$CE=ND=\dfrac{1}{2}AD=BC.$$
		Vậy $E\equiv B$ và $\dfrac{CE}{CB}=1$. 
		}{\begin{tikzpicture}[scale=1,font=\footnotesize,line join=round,line cap=round,>=stealth]
		\path (0,0) coordinate (B)++(2,0) coordinate  (C)++(1.3,1) coordinate (D)
			($(B)+(D)-(C)$)  coordinate (a)
			(barycentric cs:a=-2,D=1)  coordinate (A)
			(1,2.5)  coordinate (h)
			($(A)+(h)$)  coordinate (S)
			(barycentric cs:S=1,A=1) coordinate  (M)
			(barycentric cs:D=1,A=1)  coordinate (N);
		\draw  (S)--(A)--(B)--(C)--(D)--cycle (M)--(B)--(S)--(C);
		\draw[dashed] (A)--(D) 
			 (M)--(N)--(B);
		\foreach \p/\g in {A/150,B/240,C/-60,D/0,S/90,M/150,N/-90}  \fill[black] (\p) circle (1pt)+(\g:0.3) node{$\p$};
	\end{tikzpicture}
		 }
	} 
\end{ex}

% nguồn: Ok_TraLoiNgan_Lop11_Lan1
\begin{ex}%[1H4V4-5]%[ma:835f7945]
	Cho tứ diện $ABCD$. Gọi $E$ là điểm thuộc cạnh $AD$ sao cho $DE=2EA$, $G$ là trọng tâm tam giác $BCD$. Gọi $(\alpha)$ là mặt phẳng chứa $EG$ và song song với $BC$.  Xác định giao tuyến của mặt phẳng $(\alpha)$ và mặt phẳng $(BCD)$. Hỏi có bao nhiêu giao tuyến của mặt phẳng $(\alpha)$ với các mặt phẳng còn lại của tứ diện đi qua $A$.
	 \shortans{2}
	\loigiai{ 
		\begin{center}
			\begin{tikzpicture}[scale=1,font=\footnotesize,line join=round,line cap=round,>=stealth]
					 \def\a{4}
					\path (0:0) coordinate (B) 
						++(0:\a) coordinate  (D)
						 ++(-120:\a/2) coordinate (C)
						($(B)+(70:\a)$)  coordinate (A)
						($(A)!1/3!(D)$) coordinate  (E)
						(barycentric  cs:D=1,B=1,C=1) coordinate (G)
						($(B)!5cm!(C)$) coordinate (Bt) 
						($(G)+(Bt)-(B)$) coordinate  (Gt)
						(intersection of G--Gt  and C--D) coordinate (M)
						($(G)+(Bt)-(C)$) coordinate  (G1t)
						 (intersection of G--G1t and B--D) coordinate (N);
					 \draw[dashed] (B)--(D)
						 (E)--(G)
						 (M)--(N)--(E);
					\draw (B)--(A)--(D) 
						 (A)--(C)
						 (M)--(E)
						(B)--(C)--(D); 
					\foreach  \x/\g in {A/90,B/180,C/-45,D/0,E/45,G/-90,M/-45,N/135}
					\fill (\x) circle (1pt) 
						 ($(\g:3mm)+(\x)$) node {$\x$};
				\end{tikzpicture}
		\end{center}
		Qua $G$ đựng đường thẳng song song với $BC$ cắt $CD$ và $BD$ lần lượt tại  hai điểm $M$, $N$.\\
		Ta  có $G$ là trọng tâm $\triangle BCD \Rightarrow \dfrac{BN}{ND}=\dfrac{CM}{MD}=\dfrac{1}{2}$.\\
		 Theo đề bài, ta có $\dfrac{EA}{DE}=\dfrac{1}{2}$.\\
		Suy ra, $NE\parallel AB$ và $ME\parallel  AC \Rightarrow (ABC)\parallel (EMN)$\\
		Vậy mặt phẳng $(\alpha)$ là $(EMN)$.\\ 
		Ta có  $\heva{&N=BD\cap(EMN)\\&M=CD\cap (EMN)}\Rightarrow (EMN)\cap (BCD)=MN$.\\
		Tương tự ta tìm được giao tuyến của $(\alpha)\cap (ACD)=EM$ và  $(\alpha)\cap (ABD)=EN$.\\
		Vậy có $2$ giao tuyến của mặt phẳng $(\alpha)$ với  các mặt phẳng còn lại của tứ diện đi qua $A$.
	} 
\end{ex}


%%%%%%%%%%%%%%%% TỰ LUẬN %%%%%%%%%%%%%%%%

% nguồn: Ok_TuanLuan_Toan11_Lan2
\begin{ex}%[1H4H4-5]%[ma:fbd22634]
	Cho hình chóp $S. ABCD$ có đáy  là hình bình hành tâm $O$. Gọi $M, N$ lần lượt là trung điểm của $SA$ và $CD$
	\begin{enumerate}
		\item[a)] Chứng  minh $(OMN) \parallel(SBC)$.
		\item[b)]  Giả sử hai tam giác $SAD$ và $SAB$ là các tam giác cân tại $A$. Gọi $AE$ và $AF$ lần lượt là đường phân giác trong của hai tam giác $SAD$ và $SAB$. Chứng minh $EF\parallel(SBD)$.
	\end{enumerate}
	 \loigiai{
		\begin{center}
			\begin{tikzpicture}[scale=0.7,font=\footnotesize,line join=round,line cap=round,>=stealth]
					 \path
						(0,0) coordinate  (A)
						++(-120:3)  coordinate (D)
						 (0:6) coordinate (B)
						($(B)+(D)-(A)$) coordinate (C) 
						(80:5) coordinate (S) 
						 ($(S)!.5!(A)$) coordinate (M)
						($(C)!.5!(D)$)  coordinate (N)
						($(A)!.5!(C)$) coordinate (O) 
						($(S)!5/11!(B)$) coordinate (F) 
						 ($(S)!5/11!(D)$) coordinate (E);
					 \foreach \i in{B,C,D}{\draw (S)--(\i);};
					\draw 
						 (B)--(C)--(D);
					 \draw[dashed]
						(B)--(A)--(D) 
						 (A)--(S)
						(A)--(C) 
						(B)--(D) 
						 (A)--(E)--(F)--cycle
						(O)--(M)--(N)--cycle; 
					\foreach \i/\g in  {A/180,B/-30,C/-60,D/-120,S/90,E/-180,F/30,O/60,M/30,N/-90}
					\fill (\i) circle  (1pt)+(\g:3mm) node[scale=.7]{$\i$};
				\end{tikzpicture}
		\end{center}
		\begin{enumerate}
			\item[a)] Ta có $MO\parallel SC$ ($MO$ là đường trung  bình của tam giác $SAC$), suy ra $MO\parallel (SCB)$.\\
			Đồng thời, $NO\parallel BC$ ($NO$ là đường trung bình của tam giác  $DCB$), suy ra $NO\parallel (SCB)$.\\
			Mà $MO$, $NO\subset (OMN)$  và $MO\cap NO = O$.\\
			Vậy $(OMN)\parallel (SBC)$. 
			\item[b)] Ta có hai tam giác $SAD$ và $SAB$  là các tam giác cân tại $A$, suy ra $AE$ và $AF$ vừa là đường phân giác vừa là đường trung tuyến lần lượt của hai tam giác $SAD$ và $SAB$.\\
			 Do đó, $E$ và $F$ lần lượt là trung điểm của $SD$ và $SB$.\\
			Suy ra $EF\parallel BD$ ($EF$ là đường trung bình của tam giác $SDB$).\\ 
			Do đó,  $EF\parallel (SBD)$.
		\end{enumerate}
	 }
\end{ex}

% nguồn: Ok_TuanLuan_Toan11_Lan2
\begin{ex}%[1H4H4-5]%[ma:114c4564]
	Cho hình chóp $S.ABCD$ có đáy $ABCD$ là hình bình hành tâm $O$. Gọi $M$, $N$, $P$ lần lượt  là trung điểm của các cạnh $AB$, $CD$, $SA$. Chứng minh rằng mặt phẳng $(DMP)$ song song với mặt phẳng $(SBN)$.
	 \loigiai{
		\begin{center}
			\begin{tikzpicture}[scale=0.7,font=\footnotesize,line join=round,line cap=round,>=stealth]
					\path 
						(0,0) coordinate  (A)
						++(-120:3)  coordinate (D)
						(0:6)  coordinate (B)
						($(B)+(D)-(A)$) coordinate  (C)
						(80:5) coordinate (S) 
						($(B)!.5!(A)$)  coordinate (M)
						($(C)!.5!(D)$)  coordinate (N)
						($(S)!.5!(A)$)  coordinate (P)
						($(A)!.5!(C)$)  coordinate (O);
					\foreach  \i in{B,C,D}{\draw (S)--(\i);};
					 \draw
						 (B)--(C)--(D)
						 (S)--(N);
					\draw[dashed] 
						 (D)--(P)--(M)--cycle
						 (B)--(N)
						(S)--(A)--(C) 
						 (B)--(D)--(A)--(B);
					 \fill[purple!30,opacity=.5] (D)--(P)--(M)--cycle;
					\fill[blue!30,opacity=.5] (S)--(B)--(N)--cycle; 
					\foreach \i/\g in  {A/180,B/-30,C/-60,D/-120,S/90,O/60,M/-60,N/-90,P/30}
					\fill (\i) circle  (1pt)+(\g:3mm) node[scale=.7]{$\i$};
				\end{tikzpicture}
		\end{center}
		Vì $M$, $P$ lần lượt là trung điểm của các cạnh $AB$,  $SA$ nên $MP\parallel SB\Rightarrow MP\parallel (SBN)$.\\
		Vì $ABCD$  là hình bình hành và $M$, $N$ lần lượt là trung điểm của các cạnh $AB$, $CD$ nên $DM\parallel NB\Rightarrow DM\parallel (SBN)$.\\
		 Do đó $(DMP)\parallel (SBN)$.
	 }
\end{ex}

% nguồn: Ok_TuanLuan_Toan11_Lan2
\begin{ex}%[1H4H4-5]%[ma:2b8107e1]
	Trong không gian cho hai hình  bình hành $ABCD$ và $ABEF$ nằm trong hai mặt phẳng phân biệt. Chứng minh rằng mặt phẳng $(AFD) \parallel (BCE)$.
	 \loigiai{
		\begin{center}
			\begin{tikzpicture}[scale=0.7,font=\footnotesize,line join=round,line cap=round,>=stealth]
					 \path
						(0:0)  coordinate (A)
						(0:5) coordinate  (B)
						 (-60:3) coordinate (D)
						(60:4) coordinate  (F)
						($(B)+(D)-(A)$)  coordinate (C)
						($(B)+(F)-(A)$) coordinate (E); 
					 \draw
						(A)--(D)--(F)--cycle 
						 (C)--(E)
						 (C)--(D)
						 (E)--(F);
					\draw[dashed] 
						(A)--(B) 
						(C)--(B)--(E); 
					\foreach  \i/\g in {A/180,B/0,C/-60,D/-120,E/60,F/120}
					\fill (\i) circle (1pt)+(\g:3mm) node[scale=.7]{$\i$}; 
				\end{tikzpicture}
		\end{center}
		 Ta có
		$ 
		\heva{ 
		&AF\parallel  BE, \ AF\subset (ADF), \ BE\subset (BEC) \\
		&AD\parallel BC, \ AD\subset (ADF), \ BC\subset  (BEC)
		 }
		\Rightarrow  (ADF)\parallel (BEC)$.
	 }
\end{ex}

% nguồn: Ok_TuanLuan_Toan11_Lan2
\begin{ex}%[1H4H4-5]%[ma:1d83fbd9]
	Cho hai hình vuông $ABCD$ và $ABEF$ ở trong hai mặt phẳng phân biệt. Trên các đường chéo $AC$ và $BF$ lần lượt lấy các điểm $M, N$  sao cho $AM=BN$. Các đường thẳng song song với $AB$ vẽ từ $M$ và $N$ lần lượt cắt $AD$ và $AF$ tại $M'$ và $N'$. Chứng minh
	\begin{enumerate}
		\item[a)] $(ADF)  \parallel(BCE)$.
		\item[b)] $(DEF)  \parallel\left(MM'N'N\right)$.
	\end{enumerate}
	 \loigiai{
		\begin{center}
			\begin{tikzpicture}[scale=0.7,font=\footnotesize,line join=round,line cap=round,>=stealth]
					\path 
						(0,0) coordinate (A) 
						++(-120:3) coordinate  (D)
						++(0:5) coordinate  (C)
						($(A)+(C)-(D)$)  coordinate (B)
						($(A)+(0,5)$)  coordinate (F)
						($(F)+(C)-(D)$)  coordinate (E)
						($(A)!1/3!(D)$)  coordinate (M')
						($(A)!1/3!(C)$) coordinate (M) 
						($(A)!1/3!(F)$) coordinate (N') 
						($(B)!1/3!(F)$) coordinate  (N);
					\draw 
						(F)--(D)--(C)--(E)--cycle 
						(D)--(E)--(B)--(C); 
					\draw[dashed] 
						(C)--(A)--(B)--(F)--(A)--(D) 
						 (M')--(N')--(N)--(M)--cycle;
					 \foreach \i/\g in {A/30,B/30,C/-60,D/-150,F/120,E/60,M'/180,M/0,N/0,N'/180}
					\fill (\i) circle (1pt)+(\g:3mm) node[scale=.7]{$\i$}; 
				\end{tikzpicture}
		\end{center}
		\begin{enumerate}
			\item[a)] Ta có  $AD\parallel BC\Rightarrow AD\parallel (BCE)$; $AF\parallel BE\Rightarrow AF\parallel (BCE)$.\\
			Vì $AD\cap AF=A$ và  $AD$, $AF\subset (ADF)$ nên $(ADF) \parallel (BCE)$.
			\item[b)] Trong tam giác $ADC$, ta có $MM'\parallel CD$, theo định lí  Thales suy ra $\dfrac{AM'}{AD}=\dfrac{AM}{AC}$. \tagEX{1}
			Trong  tam giác $AFB$, ta có $NN'\parallel AB$, theo định lí Thales suy ra $\dfrac{AN'}{AF}=\dfrac{BN}{BF}$. \tagEX{2}
			Ta có $AM=BN$ (giả thiết) và $AC=BF$ (vì $ABCD$  và $ABEF$ là các hình vuông). \tagEX{3}
			Từ $(1)$,  $(2)$ và $(3)$ suy ra $\dfrac{AM'}{AD}=\dfrac{AN'}{AF}$.\\
			Theo định lí Thales đảo trong tam giác  $ADF$, suy ra $M'N' \parallel DF$.\\
			Do đó, $DF\parallel\left(MM'N'N\right)$.  \tagEX{4}
			Lại có $NN'\parallel AB$ (giả thiết),  suy ra $NN'\parallel EF$, suy ra $EF\parallel\left(MM'N'N\right)$. \tagEX{5}
			Mặt khác, $DF\cap EF=F$ và $DF$, $EF\subset(DEF)$. \tagEX{6} 
			Từ $(4)$, $(5)$ và  $(6)$ suy ra $(DEF)\parallel \left(MM'N'N\right)$.
		\end{enumerate}
	} 
\end{ex}

% nguồn: Ok_TuanLuan_Toan11_Lan2
\begin{ex}%[1H4H4-5]%[ma:c32d97dc]
	Cho hai hình bình hành $ABCD$ và $ABEF$ có tâm lần lượt là $O, O'$ và không cùng nằm trong một mặt phẳng. Gọi $M$ là trung điểm của $AB$. Chứng  minh rằng
	\begin{enumerate}
		 \item[a)] $(ADF)\parallel (BCE)$.
		\item[b)] $\left(MOO'\right)\parallel  (ADF)$.
		\item[c)] $\left(MOO'\right)\parallel (BCE)$. 
	\end{enumerate}
	\loigiai{ 
		\begin{center}
			\begin{tikzpicture}[scale=0.7,font=\footnotesize,line join=round,line cap=round,>=stealth]
					\path 
						(0,0)  coordinate (A)
						++(-135:3)  coordinate (D)
						++(0:7) coordinate (C) 
						 ($(A)+(C)-(D)$) coordinate (B)
						($(A)+(105:3)$) coordinate  (F)
						($(F)+(C)-(D)$) coordinate (E) 
						 ($(F)!1/2!(B)$) coordinate (O')
						($(A)!1/2!(B)$) coordinate (M) 
						 ($(A)!1/2!(C)$) coordinate (O);
					\draw 
						(F)--(D)--(C)--(E)--cycle 
						(E)--(B)--(C); 
					\draw[dashed] 
						 (F)--(A)--(D)
						(B)--(A) 
						 (E)--(A)--(C)
						 (D)--(B)--(F)
						 (O)--(M)--(O')--cycle;
					\foreach \i/\g in {A/180,B/-30,C/-60,D/-150,F/120,E/60,O'/90,O/-90,M/-60} 
					\fill (\i) circle (1pt)+(\g:3mm)  node[scale=.7]{$\i$};
				\end{tikzpicture}
		\end{center}
		\begin{enumerate}
			\item[a)] Ta có 
			$ 
			\heva{ 
			&AD \cap AF = \{I\} \\ 
			&AD, AF \subset  (ADF) \\
			&BC, BE \subset (BCE) \\ 
			&AD \parallel BC, \ AF \parallel  BE
			} \Rightarrow (ADF)  \parallel (BCE)
			$.\\ 
			\item[b)]  Do $O$, $O'$ lần lượt là tâm các hình bình hành nên $O$, $O'$ là trung điểm các đường chéo $AC$, $BD$ và $AE$, $BF$.\\
			Theo tính chất đường  trung bình trong tam giác, ta có
			\begin{align*}
				OO' \parallel DF,  \ OO' \parallel CE, \ OM \parallel AD, \ OM \parallel BC.
			\end{align*}
			Khi đó 
			$ 
			\heva{ 
			&OO' \cap OM \subset (MOO')  \\
			&DF,  AD \subset (ADF) \\
			&OO' \parallel DF, \ OM \parallel AD 
			} 
			\Rightarrow (MOO') \parallel (ADF)$. 
			\item[c)] Tương tự, ta có 
			 $
			 \heva{
			&OO' \cap OM \subset (MOO')  \\
			&CE, BC \subset  (BCE) \\
			&OO' \parallel CE, \ OM \parallel BC 
			} \Rightarrow (MOO')  \parallel (BCE)$.
		\end{enumerate}
	 }
\end{ex}

% nguồn: Ok_TuanLuan_Toan11_Lan2
\begin{ex}%[1H4H4-5]%[ma:63b71356]
	Cho hình chóp $S.ABCD$ có đáy $ABCD$  là hình bình hành tâm $O$, gọi $M$, $N$ lần lượt là trung điểm của $SA$, $SD$. Chứng minh $(OMN)\parallel (SBC)$.
	 \loigiai{
		\begin{center}
			\begin{tikzpicture}[scale=0.7,font=\footnotesize,line join=round,line cap=round,>=stealth]
					\path 
						(0,0)  coordinate (A)
						++(-140:3)  coordinate (D)
						++(0:4.5) coordinate (C) 
						($(A)+(C)-(D)$) coordinate  (B)
						 ($(A)!1/2!(D)$) coordinate (H)
						 ($(H)+(0,4.5)$) coordinate (S)
						($(A)!.5!(C)$) coordinate  (O)
						($(S)!.5!(A)$) coordinate (M) 
						($(S)!.5!(D)$) coordinate (N); 
					\foreach  \i in{B,C,D}{\draw (S)--(\i);};
					 \draw (B)--(C)--(D);
					 \draw[dashed]
						 (S)--(A)--(B)
						(A)--(D) 
						(A)--(C) 
						(B)--(D) 
						(O)--(M)--(N)--cycle; 
					\foreach  \i/\g in {A/150,B/-60,C/-60,D/-120,S/90,O/-90,M/30,N/150}
					\fill (\i) circle (1pt)+(\g:3mm)  node[scale=.7]{$\i$};
				\end{tikzpicture}
		\end{center}
		\noindent Ta có $MO$  là đường trung bình của tam giác $SAC$, suy ra $MO\parallel SC$. Do đó $MO\parallel(SBC)$. \tagEX{1}
		\noindent Ta có $NO$ là đường  trung bình của tam giác $SDB$, suy ra $NO\parallel SB$. Do đó $NO\parallel(SBC)$. \tagEX{2}
		\noindent Mặt khác, $NO\cap MO=O$ và  $NO$, $MO\subset (OMN)$. \tagEX{3}
		\noindent Từ $(1)$, $(2)$ và $(3)$ suy ra  $(OMN)\parallel (SBC)$.
	 }
\end{ex}

% nguồn: Ok_TuanLuan_Toan11_Lan2
\begin{ex}%[1H4H4-5]%[ma:a451cc1e]
	Cho hình chóp $S.ABCD$, gọi $G_1$, $G_2$, $G_3$ lần lượt là trọng tâm của tam giác $SAB$, $ABC$, $SAC$. Chứng  minh rằng $\left(G_1G_2G_3\right) \parallel(SBC)$.
	 \loigiai{
		\begin{center}
			\begin{tikzpicture}[scale=0.7,font=\footnotesize,line join=round,line cap=round,>=stealth]
					\path 
						(0,0) coordinate (A) 
						++(-60:3)  coordinate (B)
						++(-10:4) coordinate  (C)
						(0:7)  coordinate (D)
						(60:5) coordinate (S) 
						($(A)!.5!(B)$) coordinate (M) 
						 ($(S)!2/3!(M)$) coordinate (G_1)
						($(C)!2/3!(M)$) coordinate  (G_2)
						($(A)!.5!(C)$) coordinate  (N)
						($(S)!2/3!(N)$) coordinate  (G_3);
					\foreach \i  in {A,B,C,D}{\draw (S)--(\i);}
					\draw 
						 (A)--(B)--(C)--(D)
						 (S)--(M);
					\draw[dashed] 
						 (A)--(D)
						(G_1)--(G_2)--(G_3)--cycle 
						 (A)--(C)
						(B)--(D) 
						(C)--(M) 
						(S)--(N)--(B); 
					\foreach \i/\g in {A/120,B/-120,C/-60,D/60,S/90,G_1/150,G_2/-180,G_3/30,N/30,M/180} \fill  (\i) circle (1.3pt)+(\g:4mm) node[scale=.7]{$\i$};
				\end{tikzpicture}
		\end{center}
		Gọi $M$, $N$ lần lượt là trung điểm  của $AB$, $AC$.\\
		Khi đó  $G_1G_2$, $G_2G_3$ lần lượt là đường trung bình trong tam giác $SMC$, $SNB$.\\
		Suy  ra $G_1G_2\parallel SC$, $G_2G_3 \parallel SB$.\\
		Vì  $\heva{
		&G_1G_2,  \ G_2G_3 \subset \left(G_1G_2G_3\right)\\
		&G_1G_2\cap G_2G_3=G_2\\ 
		&G_1G_2\parallel SC \\ 
		&G_2G_3 \parallel SB 
		 }
		 \Rightarrow \left(G_1G_2G_3\right)\parallel (SBC)$.
	 }
\end{ex}

% nguồn: Ok_TuanLuan_Toan11_Lan2
\begin{ex}%[1H4H4-5]%[ma:5b2ee210]
	Cho hình hộp $ABCD.A'B'C'D'$. Hai điểm $M$, $N$ không trùng với các đầu mút, lần lượt nằm trên hai cạnh $AD$, $CC'$ sao cho $\dfrac{AM}{MD}=\dfrac{CN}{NC'}$. Xác định thiết diện của hình hộp cắt bởi mặt phẳng  qua $MN$ và song song với $\left(ACB'\right)$.
	 \loigiai{
	\begin{center}
		\begin{tikzpicture}[scale=0.7,font=\footnotesize,line join=round,line cap=round,>=stealth]
				\coordinate (A) at  (0,0);
				\coordinate (B) at  (-2,-2);
				\coordinate (C) at (4,-2); 
				\coordinate (D) at  (6,0);
				\coordinate  (A') at ($(A)+(0,5)$);
				\coordinate (B') at ($(B)+(0,5)$); 
				 \coordinate (C') at ($(C)+(0,5)$);
				\coordinate  (D') at ($(D)+(0,5)$);
				\coordinate (M)  at ($(A)!0.6!(D)$);
				\coordinate  (N) at ($(C)!0.6!(C')$);
				 \coordinate (E) at ($(C)!0.6!(D)$);
				\coordinate  (F) at ($(B')!0.6!(C')$);
				\coordinate (K) at ($(B')!0.6!(A')$); 
				 \coordinate (I) at ($(A)!0.6!(A')$);
				 \fill[yellow](K)--(F)--(N)--(E)--(M)--(I)--cycle;
				 \draw (B)--(C)--(C')--(B')--cycle (C)--(D)--(D')--(C')
					(C')--(A')--(B')--(B) 
					 (A')--(D')
					(K)--(F)--(N)--(E) 
					 (B')--(C);
				 \draw[dashed] (K)--(I)--(M)--(E)
					(B')--(A)--(C) 
					 (A')--(D)
					 (B)--(A)--(D)
					 (A')--(A);
				 \foreach \i/\g in {A/-90,B/-90,C/-45,D/0,A'/90,B'/180,C'/90,D'/0,K/180,I/180,F/90,N/0,E/0,M/-90}{\draw[fill](\i) circle (1pt)
					 ($(\i)+(\g:3mm)$) node[scale=1]{$\i$};}
			\end{tikzpicture}
	\end{center}
	Gọi $I$ là điểm trên $AA'$ sao cho $\dfrac{A I}{I A'}=\dfrac{A M}{M D}$ suy ra $I M \parallel A'D$ suy ra $I M \parallel C B'$.\\ 
	Lại có $\dfrac{A I}{I A'}=\dfrac{C N}{N C'}$ suy ra $I N  \parallel A C$ suy ra $(MNI) \parallel\left(A C B'\right)$.\\
	Do đó mặt phẳng $(MNI)$ là mặt phẳng đi qua $M, N$  và song song vưới mặt phẳng $\left(ACB'\right)$.\\
	Qua $M$ kẻ  $M E \parallel AC$; qua $N$ kẻ $N F / / B'C'$, qua $F$ kẻ $F K \parallel A'C'$.\\
	Khi  đó thiết diện cần tìm là đa giác $MENFKI$.
	 }
\end{ex}

% nguồn: Ok_TuanLuan_Toan11_Lan2
\begin{ex}%[1H4V4-5]%[ma:2b836dd9]
	Cho hình chóp $S.ABCD$ có đáy là hình thang $ABCD$, $AD\parallel BC$, $AD=2BC$. Gọi $E$, $F$, $I$ lần lượt là trung điểm của các  cạnh $SA$, $AD$, $SD$.
	\begin{enumerate}
		\item[a)] Chứng minh $(BEF)\parallel (SCD)$  và $CI\parallel (BEF)$.
		\item[b)]  Tìm giao tuyến của hai mặt phẳng $(SBC)$ và $(SAD)$.
		\item[c)] Tìm giao điểm $K$ của $FI$ với giao tuyến vừa tìm được ở  câu $b$, từ đó chứng minh $(SBF)\parallel (KCD)$.
	\end{enumerate}
	\loigiai{ 
		\begin{center}
			\begin{tikzpicture}[scale=0.7,font=\footnotesize,line join=round,line cap=round,>=stealth]
					 \path
						 (0,0) coordinate (A)
						++(-140:3) coordinate (B) 
						++(0:4)  coordinate (C)
						 (0:8) coordinate (D)
						($(A)!1/2!(B)$) coordinate  (H)
						($(H)+(0,5)$)  coordinate (S)
						($(S)!.5!(A)$)  coordinate (E)
						($(S)!.5!(D)$) coordinate (I) 
						($(A)!.5!(D)$) coordinate  (F)
						($(S)+(D)-(F)$) coordinate (K) 
						($(S)!1.5!(K)$)  coordinate (X)
						($(K)!1.5!(S)$) coordinate  (Y);
					\foreach \i in{B,C}{\draw (S)--(\i);}; 
					 \foreach \i in{A,D,F}{\draw[dashed] (S)--(\i);};
					 \draw
						 (B)--(C)--(D)
						(X)--(Y) node[below]{$d$} 
						(C)--(K)--(D); 
					\draw[dashed] 
						 (B)--(E)--(I)--(C)
						 (D)--(A)--(B)--(F)--(K)
						(E)--(F); 
					\foreach \i/\g  in {A/60,B/-120,C/-60,D/-30,S/90,E/-180,I/30,F/30,K/90}
					\fill (\i)  circle (1pt)+(\g:3mm) node[scale=.7]{$\i$};
				\end{tikzpicture}
		\end{center}
		\begin{enumerate}
			\item[a)] Ta có $EF\parallel SD$ (đường trung bình của tam giác $SAD$), suy  ra $EF\parallel (SCD)$ (vì $SD \subset (SCD)$).\\
			Ta có $BF \parallel CD$ ($BFDC$ là hình  bình hành), suy ra $BF \parallel (SCD)$ (vì $CD \subset (SCD)$).\\
			Mặt khác, $EF$, $BF\subset (BEF)$  và $EF\cap BF = F$.\\
			Vậy $(BEF) \parallel (SCD)$.\\ 
			Ta có $EI\parallel BC$ (vì cùng song song với $AD$) và $EI=BC=\dfrac{1}{2}AD$  (vì $EI$ là đường trung bình của tam giác $SAD$ và $AD=2BC$).\\
			Tứ giác $EICB$ là hình bình hành, nên $CI\parallel  BE$.\\
			 Mặt khác, $BE\subset (BEF)$ suy ra $CI\parallel (BEF)$.
			\item[b)] Ta  có
			 $\heva{
			 &S\in (SBC)\cap (SAD)\\
			 &BC \parallel AD\\
			&BC  \subset (SBC)\\
			&AD\subset (SAD). 
			 }$\\
			Suy ra giao tuyến của $(SBC)$ và $(SAD)$ là đường thẳng $d$ đi qua $S$ và $d\parallel BC\parallel AD$. 
			\item[c)] Xét tam giác  $SAD$ có $I$ là trung điểm của $SD$, $F$ là trung điểm của $AD$.\\
			Suy ra $IF$ là đường trung bình của tam giác $SAD$, suy ra $IF\parallel  SA$. Do đó, $KF \parallel SA$. \tagEX{1}
			Mặt khác, $SK \parallel  AF$. \tagEX{2}
			Từ $(1)$ và $(2)$ suy ra $SKFA$ là hình  bình hành.\\
			Suy ra $SK=AF$. Lại có $SK=FD$ (vì $AF = FD$).\\ 
			Tứ giác  $SKDF$ có $SK=FD$ và $SK\parallel FD$, nên $SKDF$ là hình bình hành.\\
			Suy ra $SF \parallel  KD$.\\
			Ta có $SF \parallel KD$,  $BF \parallel DC$, $F = SF \cap BF$, $D = KD \cap CD$.\\
			Suy ra $(SBF) \parallel (KCD)$. 
		\end{enumerate}
	 }
\end{ex}

% nguồn: Ok_TuanLuan_Toan11_Lan2
\begin{ex}%[1H4V4-5]%[ma:f5de99ea]
	Cho hình chóp $S.ABCD$ có đáy $ABCD$ là hình bình hành và $M$, $N$ lần lượt là  trung điểm của $AB$, $CD$. Gọi $(P)$ là mặt phẳng đi qua $MN$ và song song với mặt phẳng $(SAD)$. Tìm giao tuyến của các mặt của hình chóp với mặt phẳng $(P)$.
	\loigiai{ 
		\begin{center}
			\begin{tikzpicture}[scale=1.5,font=\footnotesize,line join=round,line cap=round,>=stealth]
					\def\a{3} 
					 \def\b{1.5}
					 \def\g{30}
					 \def\h{2}
					\path 
						(0:0)  coordinate (A)--++(\g:\b) coordinate (B)--++(0:\a) coordinate (C)--++(\g-180:\b) coordinate (D)
						 (A)--++(70:3.5) coordinate (S);
					\draw ($(A)!.5!(B)$) coordinate  (M);
					\draw ($(C)!.5!(D)$) coordinate (N); 
					\draw ($(S)!.5!(B)$) coordinate (F); 
					\draw ($(S)!.5!(C)$) coordinate (E); 
					\foreach \x/\g  in {A/180,B/150,C/0,D/-30,S/90,M/160,N/-60,F/-50,E/40}
					\fill(\x)  circle (1pt)
						($(\x)+(\g:3mm)$)  node{$\x$};
					 \draw[dashed] (S)--(B)--(A)
						(C)--(B) 
						 (N)--(M)--(F)--(E);
					\draw  (S)--(A)--(D)--(C)--(S)--(D)
						(E)--(N); 
				\end{tikzpicture}
		\end{center}
		$(P) \cap (ABCD)=MN$ với $MN  \parallel BC \parallel AD$.\\
		$(P) \cap (SAB)=MF$ với $MF \parallel SA  (F\in SB)$.\\
		$(P) \cap (SDC)=N  E$ với $NE \parallel SD (E \in SC)$.\\
		$(P) \cap (SBC)=EF$ với  $EF \parallel BC \parallel MN$.
	 }
\end{ex}
